% GENERATED FILE. Edit canonical lessons, then run npm run generate:books.
% canonical-source-hash: fnv1a64:dd01cdcf0b3ed376
% canonical-lessons: GU-C291-apaharan, GU-C291-gum, GU-C291-ghervu, GU-C291-kedi, GU-C291-chodavvu

\chapter{Surrounded}
\label{ch:gu-surrounded}
\hlchaptermodality{291}

\begin{chapteropening}
\textbf{By the end of this chapter:} \emph{I can tell how the police surrounded a house and freed a captive.}

Complete the last lesson of chapter 291: I can tell how the police surrounded a house and freed a captive.
\end{chapteropening}

\section[\texorpdfstring{\gu{અપહરણ} (apaharaṇ)}{અપહરણ (apaharaṇ)}]{\gu{અપહરણ} (apaharaṇ) — a kidnapping}

\label{lesson:GU-C291-apaharan}



\begin{quote}
Before the new one: say the Gujarati for to threaten, then the Gujarati for to snatch.
\end{quote}

\subsection*{You'll want to know: \gu{અપહરણ}}
\textbf{\gu{અપહરણ}} — \emph{apaharaṇ} — \textquotedblleft{}a kidnapping\textquotedblright{}.

\begin{quote}
\gu{ગઈકાલે} \gu{એક} \gu{વેપારીના} \gu{દીકરાનું} \gu{અપહરણ} \gu{થયું}.
\end{quote}

\begin{quote}
\emph{gaīkāle ek vepārīnā dīkrānũ apaharaṇ thayũ.}
\end{quote}

\begin{quote}
Yesterday a merchant's son was kidnapped.
\end{quote}

A neuter noun.

\begin{cousinweb}
\textbf{\gu{અપહરણ}} (\emph{apaharaṇ}) is a tatsama, Sanskrit \emph{apaharaṇa}, \textquotedblleft{}a carrying away, stealing\textquotedblright{} (Monier-Williams). In Gujarati it is the carrying off of a person: \textbf{\gu{દીકરાનું} \gu{અપહરણ} \gu{થયું}}.
\end{cousinweb}

\begin{grammarlens}[title={What you've built}]
One more word for telling what happened.
\end{grammarlens}

\subsection*{Your turn}
\begin{itemize}
\raggedright
  \item \emph{Say it:} \emph{apaharaṇ}
  \item \emph{Say it:} \emph{apaharaṇ}, once more
  \item \emph{Say it:} \emph{gaīkāle ek vepārīnā dīkrānũ apaharaṇ thayũ.} — Yesterday a merchant's son was kidnapped.
  \item \emph{Recall:} say the Gujarati for to threaten, then the Gujarati for to snatch, then say \emph{apaharaṇ} again
\end{itemize}

\begin{tcolorbox}[breakable,colback=teal!4,colframe=teal!35!black,title={Before you move on}]
What does \gu{અપહરણ} mean? (\textquotedblleft{}A kidnapping\textquotedblright{}.) Say it once more.
\end{tcolorbox}
\section[\texorpdfstring{\gu{ગુમ} (gum)}{ગુમ (gum)}]{\gu{ગુમ} (gum) — missing}

\label{lesson:GU-C291-gum}



\begin{quote}
Before the new one: say the Gujarati for to snatch, then the Gujarati for a kidnapping.
\end{quote}

\subsection*{You'll want to know: \gu{ગુમ}}
\textbf{\gu{ગુમ}} — \emph{gum} — \textquotedblleft{}missing\textquotedblright{}.

\begin{quote}
\gu{છોકરો} \gu{ત્રણ} \gu{દિવસથી} \gu{ગુમ} \gu{હતો}.
\end{quote}

\begin{quote}
\emph{chokro traṇ divasthī gum hato.}
\end{quote}

\begin{quote}
The boy had been missing for three days.
\end{quote}

\begin{cousinweb}
\textbf{\gu{ગુમ}} (\emph{gum}) is usually traced to Persian \emph{gum}, \textquotedblleft{}lost\textquotedblright{}. It describes a person or a thing nobody can find: \textbf{\gu{છોકરો} \gu{ગુમ} \gu{છે}}; \textbf{\gu{ગુમ} \gu{થયો}}, went missing.
\end{cousinweb}

\begin{grammarlens}[title={What you've built}]
One more word for telling what happened.
\end{grammarlens}

\subsection*{Your turn}
\begin{itemize}
\raggedright
  \item \emph{Say it:} \emph{gum}
  \item \emph{Say it:} \emph{gum}, once more
  \item \emph{Say it:} \emph{chokro traṇ divasthī gum hato.} — The boy had been missing for three days.
  \item \emph{Recall:} say the Gujarati for to snatch, then the Gujarati for a kidnapping, then say \emph{gum} again
\end{itemize}

\begin{tcolorbox}[breakable,colback=teal!4,colframe=teal!35!black,title={Before you move on}]
What does \gu{ગુમ} mean? (\textquotedblleft{}Missing\textquotedblright{}.) Say it once more.
\end{tcolorbox}
\section[\texorpdfstring{\gu{ઘેરવું} (ghervũ)}{ઘેરવું (ghervũ)}]{\gu{ઘેરવું} (ghervũ) — to surround}

\label{lesson:GU-C291-ghervu}



\begin{quote}
Before the new one: say the Gujarati for a kidnapping, then the Gujarati for missing.
\end{quote}

\subsection*{You'll want to know: \gu{ઘેરવું}}
\textbf{\gu{ઘેરવું}} — \emph{ghervũ} — \textquotedblleft{}to surround\textquotedblright{}.

\begin{quote}
\gu{પોલીસે} \gu{આખું} \gu{મકાન} \gu{ઘેર્યું}.
\end{quote}

\begin{quote}
\emph{polīse ākhũ makān gheryũ.}
\end{quote}

\begin{quote}
The police surrounded the whole building.
\end{quote}

Past, with \gu{તેણે} (he or she …): \gu{તેણે} \gu{મકાન} \gu{ઘેર્યું}, \emph{teṇe makān gheryũ}.

\begin{cousinweb}
\textbf{\gu{ઘેરવું}} (\emph{ghervũ}) comes from \emph{*ghēr-}, \textquotedblleft{}to make go round\textquotedblright{}, a form CDIAL reconstructs, and CDIAL lists the Gujarati verb, \textquotedblleft{}to surround\textquotedblright{} (CDIAL 4474).
\end{cousinweb}

\begin{grammarlens}[title={What you've built}]
One more word for telling what happened.
\end{grammarlens}

\subsection*{Your turn}
\begin{itemize}
\raggedright
  \item \emph{Say it:} \emph{ghervũ}
  \item \emph{Say it:} \emph{ghervũ}, once more
  \item \emph{Say it:} \emph{polīse ākhũ makān gheryũ.} — The police surrounded the whole building.
  \item \emph{Recall:} say the Gujarati for a kidnapping, then the Gujarati for missing, then say \emph{ghervũ} again
\end{itemize}

\begin{tcolorbox}[breakable,colback=teal!4,colframe=teal!35!black,title={Before you move on}]
What does \gu{ઘેરવું} mean? (\textquotedblleft{}To surround\textquotedblright{}.) Say it once more.
\end{tcolorbox}
\section[\texorpdfstring{\gu{કેદી} (kedī)}{કેદી (kedī)}]{\gu{કેદી} (kedī) — a prisoner; a captive}

\label{lesson:GU-C291-kedi}



\begin{quote}
Before the new one: say the Gujarati for missing, then the Gujarati for to surround.
\end{quote}

\subsection*{You'll want to know: \gu{કેદી}}
\textbf{\gu{કેદી}} — \emph{kedī} — \textquotedblleft{}a prisoner; a captive\textquotedblright{}.

\begin{quote}
\gu{જો} \gu{પોલીસ} \gu{સમયસર} \gu{પહોંચી} \gu{ગઈ} \gu{તો} \gu{કેદી} \gu{બચી} \gu{જશે}.
\end{quote}

\begin{quote}
\emph{jo polīs samaysar pahõcī gaī to kedī bacī jaśe.}
\end{quote}

\begin{quote}
If the police get there in time, the captive will be saved.
\end{quote}

A masculine or feminine noun, by the person meant.

\begin{cousinweb}
\textbf{\gu{કેદી}} (\emph{kedī}) is usually traced to Persian \emph{qaidī}, \textquotedblleft{}a prisoner\textquotedblright{}, from \emph{qaid}, captivity, a word Persian took from Arabic. It is used of a man or a woman.
\end{cousinweb}

\begin{grammarlens}[title={What you've built}]
One more word for telling what happened.
\end{grammarlens}

\subsection*{Your turn}
\begin{itemize}
\raggedright
  \item \emph{Say it:} \emph{kedī}
  \item \emph{Say it:} \emph{kedī}, once more
  \item \emph{Say it:} \emph{jo polīs samaysar pahõcī gaī to kedī bacī jaśe.} — If the police get there in time, the captive will be saved.
  \item \emph{Recall:} say the Gujarati for missing, then the Gujarati for to surround, then say \emph{kedī} again
\end{itemize}

\begin{tcolorbox}[breakable,colback=teal!4,colframe=teal!35!black,title={Before you move on}]
What does \gu{કેદી} mean? (\textquotedblleft{}A prisoner; a captive\textquotedblright{}.) Say it once more.
\end{tcolorbox}
\section[\texorpdfstring{\gu{છોડાવવું} (choḍāvvũ)}{છોડાવવું (choḍāvvũ)}]{\gu{છોડાવવું} (choḍāvvũ) — to set free; to get released}

\label{lesson:GU-C291-chodavvu}



\begin{quote}
Before the new one: say the Gujarati for to surround, then the Gujarati for a prisoner.
\end{quote}

\subsection*{You'll want to know: \gu{છોડાવવું}}
\textbf{\gu{છોડાવવું}} — \emph{choḍāvvũ} — \textquotedblleft{}to set free; to get released\textquotedblright{}.

\begin{quote}
\gu{રાત્રે} \gu{પોલીસે} \gu{છોકરાને} \gu{છોડાવ્યો}.
\end{quote}

\begin{quote}
\emph{rātre polīse chokrāne choḍāvyo.}
\end{quote}

\begin{quote}
During the night the police freed the boy.
\end{quote}

Past, with \gu{તેણે} (he or she …): \gu{તેણે} \gu{છોકરાને} \gu{છોડાવ્યો}, \emph{teṇe chokrāne choḍāvyo}.

\begin{cousinweb}
\textbf{\gu{છોડાવવું}} (\emph{choḍāvvũ}) is the causative of \textbf{\gu{છોડવું}}, to let go, which CDIAL sets under Sanskrit \emph{kṣōṭayati} (CDIAL 3747): to make someone be let go, and so to get a captive freed.
\end{cousinweb}

\begin{grammarlens}[title={What you've built}]
One more word for telling what happened.
\end{grammarlens}

\subsection*{Your turn}
\begin{itemize}
\raggedright
  \item \emph{Say it:} \emph{choḍāvvũ}
  \item \emph{Say it:} \emph{choḍāvvũ}, once more
  \item \emph{Say it:} \emph{rātre polīse chokrāne choḍāvyo.} — During the night the police freed the boy.
  \item \emph{Recall:} say the Gujarati for to surround, then the Gujarati for a prisoner, then say \emph{choḍāvvũ} again
\end{itemize}

\begin{tcolorbox}[breakable,colback=teal!4,colframe=teal!35!black,title={Before you move on}]
What does \gu{છોડાવવું} mean? (\textquotedblleft{}To set free; to get released\textquotedblright{}.) Say it once more.
\end{tcolorbox}
